\documentclass[11pt]{letter}
\usepackage[a4paper,margin=2.5cm]{geometry}
\usepackage[T1]{fontenc}
\usepackage{lmodern}
\usepackage{microtype}
\signature{Mohamed Karray\\Corresponding author, on behalf of all authors\\ESME Research Lab, Ivry-sur-Seine, France\\mohamed.karray@esme.fr}
\address{ESME Research Lab\\Ivry-sur-Seine, France}
\date{\today}

\begin{document}
\begin{letter}{The Editor-in-Chief\\Journal of Intelligent \& Robotic Systems}

\opening{Dear Editor,}

We submit our manuscript \emph{Digital-Twin Pretraining for Label-Efficient
Skeleton Action Recognition in Industrial Workcells} for consideration as an
original research article.

Collaborative robots must recognise what workers are doing, yet labelling
physical recordings for every workcell is costly. Digital twins label actions
automatically, but whether a recogniser trained on them transfers to real
workers has not been tested rigorously. Using the public InHARD-DT and InHARD
datasets with subject-disjoint folds and three seeds, we show that unsupervised
adaptation does not close this gap and explain why. We then propose a
label-efficient strategy with which a twin-initialised model matches
physical-only training while using five to twenty times fewer labelled clips.

The work fits the journal because it addresses robot perception for
human-robot collaboration in manufacturing and quantifies how much annotation
a digital twin can save in practice.

The manuscript is original, is not under consideration elsewhere, and has
been approved by all authors, who declare no competing interests.

\closing{Yours sincerely,}

\end{letter}
\end{document}
